\section{Product covers and semilinear factor actions}\label{sec:blocks}

We now prepare the klt case of the curve weight bound.  A finite
cover decomposes a klt log Calabi--Yau pair into factors with controlled
holomorphic forms.  The cover need not be Galois, however, and its
Galois closure need not be a product.  The purpose of this section is
to retain the factor directions on that closure and to show that a
bounded power of each finite transformation acts regularly on suitable
finite covers of the individual factors.  This is the comparison
construction of \cite[Lemma~3.5]{U}, extended to a rationally connected
factor carrying the boundary.

\subsection{The factors and their automorphisms}

A finite surjective morphism between normal varieties is
\emph{quasi-\'etale} if it is \'etale in codimension one.  Over an
algebraically closed field of characteristic zero, purity implies that
such a morphism is \'etale over the smooth locus of its target.
For a normal variety $F$, we write $\Omega_F^{[j]}$ for the reflexive
extension of $j$-forms from its smooth locus.

\begin{proposition}[Product decomposition]\label{prop:product-cover}
Let $(V,B)$ be a projective klt pair over $\C$ with $B\geq0$,
$K_V$ rational Cartier, and $K_V+B\simq0$.  There is a finite
quasi-\'etale cover $\pi:P\to V$ and a decomposition
\[
 (P,\pi^*B)
 \simeq (F,H)\times A\times\prod_j Y_j\times\prod_k Z_k,
\]
where $(F,H)$ is a rationally connected klt log Calabi--Yau pair,
$A$ is an abelian variety, and the $Y_j$ and $Z_k$ are respectively
irreducible Calabi--Yau and irreducible holomorphic symplectic
varieties in the singular Beauville--Bogomolov decomposition.
All factors are $\Q$-Gorenstein.  Point factors are omitted, and
the boundary is pulled back entirely from $F$.
\end{proposition}

\begin{proof}
The product decomposition, including its assertion about the boundary,
is \cite[Theorem~1.3]{MW}.  Its boundary-free factors have the stated
properties by the singular Beauville--Bogomolov theorem
\cite[Definition~1.4 and Theorem~1.5]{HP}.  Since $K_V$ is rational
Cartier, quasi-\'etaleness makes $K_P$ rational Cartier.  Restricting
the canonical sheaf in product charts, with the complementary points
smooth, gives the same property on every factor.  Restriction of
$K_P+\pi^*B\simq0$ gives $K_F+H\simq0$.
\end{proof}

We use two properties of the nonabelian boundary-free factors.  They
and all their connected normal finite quasi-\'etale covers have
canonical singularities and Cartier trivial canonical divisor.  On
each such cover the reflexive form algebra is generated by a top form
in the Calabi--Yau case and by the pulled-back symplectic form in the
symplectic case \cite[Definition~1.4]{HP}.  In particular there are no
reflexive one-forms or vector fields: contraction with a volume form,
or with the symplectic form, proves the latter assertion.  Connected
normal finite quasi-\'etale covers of an abelian variety are \'etale
by purity and are again abelian varieties after an origin is chosen.

For the rationally connected factor, vector fields must be required
to preserve the boundary.  If $D$ is a reduced Weil divisor on a
normal variety, $\mathcal T_F(-\log D)$ denotes the subsheaf of
derivations preserving the reduced ideal of every component of $D$.
This condition can be checked at height-one primes and then extended
reflexively.

\begin{lemma}[Automorphisms of a rationally connected pair]
\label{lem:rc-aut}
Let $(F,H)$ be a rationally connected projective klt pair over $\C$
with $H\geq0$ and $K_F+H\simq0$.  Then
\[
 H^j(F,\OO_F)=0\quad(j>0),\qquad
 H^0(F,\Omega_F^{[1]})=0,\qquad
 H^0\bigl(F,\mathcal T_F(-\log\Supp H)\bigr)=0.
\]
For every ample line bundle $L$, the group of automorphisms of
$(F,H)$ preserving $L$ is finite.
\end{lemma}

\begin{proof}
A smooth projective resolution of $F$ is rationally connected; one
may lift rational curves through general smooth points, or use the
rational connectedness results for klt varieties in \cite{HM}.
Its positive-degree structure-sheaf cohomology vanishes.  Klt
singularities are rational, so the same is true on $F$.  The
extension theorem for klt differential forms \cite{GKKP} then gives
the asserted vanishing of reflexive one-forms.

Choose $a>0$ and a rational $a$-canonical form $\theta$ with
\[
                         \Div(\theta)+aH=0.
\]
On the smooth locus, the density $|\theta|^{2/a}$ defines a positive
measure.  Pull it to a log resolution.  Every divisorial order of
the pulled-back form is strictly greater than $-a$, because the pair
is klt.  In SNC coordinates this is exactly the local integrability
condition for the density.  Thus it defines a finite nonzero measure
$\mu$ on $F$, with no mass on proper algebraic subsets.

An automorphism preserving $H$ multiplies $\theta$ by a nonzero
constant.  Its pullback therefore multiplies $\mu$ by a positive
constant; comparison of total masses makes that constant one.
Consequently every automorphism of the pair preserves $\mu$.

There can be no nontrivial additive or multiplicative algebraic
one-parameter subgroup of such automorphisms.  Indeed, iterate the
element $1\in\mathbb G_a$ or $2\in\mathbb G_m$ on a nonfixed
point.  The orbit map extends across infinity by projectivity, so
these iterates converge to a subgroup-fixed point.  The nonfixed
locus has full $\mu$-measure.  Such convergence contradicts
Poincar\'e recurrence for the resulting invertible transformation of
the finite measure space $(F,\mu)$.

The group preserving $L$ and $H$ is a linear algebraic group: a
sufficiently high power of $L$ realizes it as a closed subgroup of a
projective linear group.  Every positive-dimensional linear
algebraic group over $\C$ contains a copy of $\mathbb G_a$ or
$\mathbb G_m$, so this group is finite.  Moreover, the identity
component of the automorphism group preserving $H$ preserves $L$.
Indeed its variation of $L$ lies in $\operatorname{Pic}^0(F)$,
which is zero because $H^1(F,\OO_F)=0$.  Its Lie algebra therefore
vanishes.  In characteristic zero this Lie algebra is precisely the
space of vector fields tangent to $\Supp H$: a connected group
preserving the support fixes its components and their coefficients.
This proves the last vanishing.
\end{proof}

\begin{lemma}[A choice stable under further covers]\label{lem:minimal-rc}
The cover in Proposition~\ref{prop:product-cover} can be chosen so
that every connected normal finite quasi-\'etale cover of $F$ is
rationally connected.  On every such cover, with the pulled-back
boundary, the conclusions of Lemma~\ref{lem:rc-aut} hold.
\end{lemma}

\begin{proof}
Among all covers in Proposition~\ref{prop:product-cover}, choose
one with $\dim F$ minimal, taking this dimension to be zero when
the factor is absent.  Let $F'\to F$ be a connected normal finite
quasi-\'etale cover.  The pulled-back pair is klt and log
Calabi--Yau, and $K_{F'}$ is rational Cartier.  Apply
Proposition~\ref{prop:product-cover} to it.  If $F'$ were not
rationally connected, the resulting product would have a
positive-dimensional boundary-free part: otherwise rational
connectedness would descend from its rationally connected factor
along a finite surjection.  Replacing $F$ by this product would
therefore give a product cover of $V$ with a smaller rationally
connected factor, a contradiction.  Lemma~\ref{lem:rc-aut} now
applies to $(F',H')$.
\end{proof}

\subsection{Separating finite actions on a comparison cover}

We refer to the positive-dimensional factors just described as
\emph{blocks}, collecting the whole abelian part into one block.
The rationally connected block is always chosen as in
Lemma~\ref{lem:minimal-rc}.  For a block over a nonclosed field,
the corresponding properties under connected normal finite
quasi-\'etale covers, including the form and vector-field
vanishings above, are imposed after algebraic closure.  In the
application the field is a finite extension of a
complex curve function field.  Its algebraic closure is abstractly
isomorphic to $\C$, so the preceding complex projective results
apply, and their finite algebraic data descend to a finite extension.

A \emph{semilinear automorphism} of a variety over a field $k$ is an
automorphism together with an automorphism of $k$ over which it lies;
equivalently, it is a $k$-isomorphism to the corresponding field
conjugate.  Differentials below are relative to $k$.

\begin{proposition}[Semilinear comparison with the blocks]
\label{prop:semilinear}
Let $k$ be a characteristic-zero field, and let
$P=\prod_{i\in I}P_i$ be a product of geometrically integral
projective blocks as above.  Let $H_P$ be the pullback of the
boundary on its rationally connected block, or zero if that block
is absent.  Suppose
\[
                         \rho:R\longrightarrow P
\]
is a finite quasi-\'etale morphism from a geometrically integral
normal variety, and a finite group $G$ acts semilinearly on
$(R,H_R)$, where $H_R=\rho^*H_P$.  Put $s=\dim R>0$.
There are normal geometrically integral varieties $R_i$, finite
quasi-\'etale maps $R_i\to P_i$, and a finite quasi-\'etale map
$R\to\prod_iR_i$ such that every $g^{s!}$, for $g\in G$, induces
regular semilinear automorphisms of the pairs $(R_i,H_i)$.
Here $H_i$ is the pulled-back boundary on the rationally connected
block and is zero on the other blocks.  The projections from $R$
are equivariant for these automorphisms.  Each geometric $R_i$
has the same block properties as $P_i$.
\end{proposition}

\begin{proof}
We first show that a bounded power preserves the factor directions on $R$.
We then recover their constant fields and show that the resulting
birational actions on the finite factor covers are regular.

\smallskip\noindent\emph{The factor directions.}
On a smooth big open where $\rho$ is \'etale, the product tangent
directions pull back to a splitting.  Reflexive extension gives
\[
                    \mathcal T_R=\bigoplus_{i\in I}\mathcal E_i.
\]
Consider the finite-dimensional $k$-algebra
\[
 \mathcal A=
 \{u\in\operatorname{End}_R(\mathcal T_R):
 u(\mathcal T_R(-\log\Supp H_R))
       \subseteq\mathcal T_R(-\log\Supp H_R)\}.
\]
Every projection onto $\mathcal E_i$ belongs to $\mathcal A$.
We claim that every member of $\mathcal A$ has zero entries
between distinct summands.

This may be checked after extending $k$ to an algebraic closure.
Fix general smooth points in all but one factor, say $P_i$, and
take the corresponding slice of $R$.  After shrinking the space of
complementary points, its connected components are normal and
finite quasi-\'etale over $P_i$.  To see this, the generic slice is
normal by localization and geometrically normal in characteristic
zero.  The projection to the complementary factors is projective,
so generic flatness and openness of geometric normality give
normality of the entire fibres after shrinking that base.  Every
component has dimension $\dim P_i$:
the fibre-dimension inequality gives the lower bound, and finiteness
over $P_i$ gives the upper bound.  Its finite image is therefore all
of $P_i$.  The closed subset omitted from the \'etale product
charts has codimension at least two on a general slice: its
components not dominating the complementary factors can be avoided,
and the others have fibre dimension at most $\dim P_i-2$.
On its smooth big open, $\mathcal E_i$ is the tangent sheaf of
the slice, and every complementary summand is trivial.

If either of two distinct summands is nonabelian and boundary-free,
slice in that direction.  A homomorphism in either direction gives
reflexive one-forms or vector fields on a finite quasi-\'etale
cover of that block, and hence vanishes.  The only remaining
possibility is an abelian and a rationally connected summand.
Slice in the rationally connected direction.  An entry from this
summand to the abelian one gives reflexive one-forms.  An entry in
the reverse direction gives vector fields tangent to the pulled-back
boundary, because the endomorphism preserves the log tangent
subsheaf.  Both vanish by Lemmas~\ref{lem:rc-aut} and
\ref{lem:minimal-rc}.  These slices cover a dense open of $R$,
which proves the claim.

The block projections are therefore central idempotents of
$\mathcal A$.  Its primitive central idempotents give nonzero
direct summands of $\mathcal T_R$, each of positive generic rank,
so there are at most $s$ of them.  Conjugation by $G$ acts on
$\mathcal A$, since $G$ preserves the pair; it is a ring
automorphism even for a semilinear action.  It permutes these
primitive central idempotents.  Thus $g^{s!}$ fixes every one of
them and consequently every block projection.  This argument takes
place over $k$ itself and requires no extension of $G$ to an
algebraic closure.

\smallskip\noindent\emph{The finite factor covers.}
Let $k(R_i)$ be the relative algebraic closure of $k(P_i)$ in
$k(R)$, and let
\[
                         R\longrightarrow R_i\longrightarrow P_i
\]
be the Stein factorization of the projection.  At the generic point
of $R$, the complementary directions span
$\operatorname{Der}_{k(P_i)}k(R)$.  In characteristic zero their
common constants are exactly $k(R_i)$.  Preservation of the
directions therefore gives semilinear birational transformations of
$R_i$ under every $g^{s!}$.

The fields $k(R_i)$ are regular over $k$, since they lie in the
regular extension $k(R)/k$; hence $R_i$ is geometrically integral.
The extension obtained by adjoining the complementary product
factors to $k(P_i)$ is regular and linearly disjoint from
$k(R_i)/k(P_i)$.  Consequently the normal product
$R_i\times\prod_{j\ne i}P_j$ is an intermediate finite cover of
$P$.  Ramification over a prime of $P_i$ would persist on this
product and on a prolongation to $R$, contradicting
quasi-\'etaleness of $R\to P$.  This proves that $R_i\to P_i$
is quasi-\'etale.

The product map $R\to\prod_iR_i$ is proper with finite fibres,
because its composite to $P$ is finite.  It is therefore finite.
Its image has the dimension of the integral target, so it is
surjective, and the same ramification argument makes it
quasi-\'etale.  In particular every fibre of $R\to R_i$ has
dimension $s-\dim R_i$.

\smallskip\noindent\emph{Regularity of the factor actions.}
Fix a transformation preserving the factor directions.  The two
morphisms from $R$ to $R_i$---the projection and its composition
with the transformation---map $R$ onto the closed graph
$\Gamma_i$ of the induced birational factor transformation.
Their target is replaced by its field conjugate in the semilinear
case.  Every fibre of $R\to\Gamma_i$ has dimension at least
$s-\dim R_i$.  A positive-dimensional fibre of either graph
projection would therefore give a fibre of $R\to R_i$ of larger
dimension.  Both graph projections are thus finite and birational;
normality of their targets makes them isomorphisms.  This proves
regularity.

Finally, equality of divisors on $\prod_iR_i$ can be checked after
pullback along the finite surjection from $R$.  Since its entire
boundary is pulled back from the rationally connected factor,
invariance of $H_R$ implies invariance of that factor's boundary.
Hence the factor actions preserve their boundaries.  The asserted
geometric properties of the $R_i$ follow
from their being quasi-\'etale covers of the chosen blocks.
\end{proof}

After the bounded power, the action on $\prod_iR_i$ is the product
of its regular factor actions.  The factor degrees and the orders
of these automorphisms may be unbounded.

\subsection{Block forms over a curve field}

We now collect the exact data needed for degeneration, including
the degrees of the log pluriforms on the factors.

\begin{corollary}[Comparison cover and block forms]\label{cor:block-forms}
Let $K=\C(C)$ for a smooth projective curve $C$.  Let $(V,B)$
be a geometrically integral normal projective $s$-fold pair over
$K$, with $s>0$, geometrically klt, $B\geq0$, and $K_V$
rational Cartier.  Suppose that, for an integer $m>0$, a rational $m$-canonical form
$\phi$ satisfies $\Div(\phi)+mB=0$.

There is a finite Galois extension $K_1/K$, a geometrically integral
normal projective variety $R/K_1$, and finite quasi-\'etale maps
\[
\begin{tikzcd}[column sep=large]
 R \arrow[r] & \displaystyle\prod_{i\in I}R_i
 \arrow[r] & \displaystyle\prod_{i\in I}P_i
 \arrow[r] & V_{K_1}
\end{tikzcd}
\]
with the following properties.  The field $K_1(R)$ is finite Galois
over $K(V)$; its Galois group acts semilinearly on $R$, maps onto
$\operatorname{Gal}(K_1/K)$, and preserves the pulled-back pair.
For each element $g$ of this group, $g^{s!}$ induces regular
semilinear actions on the pairs $(R_i,H_i)$.  There is at most one
rationally connected block, chosen with the cover-stability of
Lemma~\ref{lem:minimal-rc}, and all other blocks are boundary-free
of the types in Proposition~\ref{prop:product-cover}.

Put $p_i=m$ on the rationally connected block and $p_i=1$ on the
other blocks.  There are rational $p_i$-canonical forms $\eta_i$
over $K_1$ satisfying
\begin{equation}\label{eq:block-trivializations}
                       \Div(\eta_i)+p_iH_i=0.
\end{equation}
On $R$, after pulling back all forms, one has
\begin{equation}\label{eq:block-product-form}
                 \phi=a\bigwedge_{i\in I}\eta_i^{\otimes m/p_i}
                 \qquad\text{for some }a\in K_1^*.
\end{equation}
The wedge uses any fixed ordering of the blocks and the natural
identification of their relative canonical lines.
\end{corollary}

\begin{proof}
Apply Proposition~\ref{prop:product-cover} and
Lemma~\ref{lem:minimal-rc} to the geometric generic pair, and
define the resulting product cover over a finite extension of $K$.
Take a Galois closure $L$ of its total function field over the
original field $K(V)$, and let $K_1$ be the relative algebraic
closure of $K$ in $L$.  The extension $K(V)/K$ is regular.
It follows that $K_1/K$ is finite Galois and that restriction maps
$\operatorname{Gal}(L/K(V))$ onto
$\operatorname{Gal}(K_1/K)$.  Let $R$ be the normalization of
$V_{K_1}$ in $L$.  Since $K_1$ is algebraically closed in $L$,
$R$ is geometrically integral.

Over an algebraic closure of $K$, the conjugate product covers are
all \'etale over the smooth locus of the original variety.  Their
compositum, and hence the Galois closure, has the same property.
Thus $R$ is finite quasi-\'etale over the product and over
$V_{K_1}$.  Its Galois transformations preserve the crepant
pullback of $B$.  Proposition~\ref{prop:semilinear} now constructs
the $R_i$ and supplies the factor actions.

The boundary-free blocks and their covers have Cartier trivial
canonical divisor over the algebraic closure, so they possess
degree-one volume forms there.  Still over the algebraic closure,
on the rationally connected factor,
restrict the pulled-back degree-$m$ trivialization from $V$ to
the product factor, choosing the complementary points smooth.
This gives a degree-$m$ log trivialization; pull it back to its
Stein block $R_i$.

Here $p_iH_i$ is integral: on the rationally connected block this
follows from the integrality of $mB$, quasi-\'etale pullback, and
restriction to the product factor; on the other blocks $H_i=0$.
Thus $p_iK_{R_i}+p_iH_i$ is an integral Weil divisor whose
divisorial sheaf becomes trivial after algebraic closure.
Lemma~\ref{lem:trivialization-descent} descends this trivialization
in degree $p_i$, giving \eqref{eq:block-trivializations}.  The product
in \eqref{eq:block-product-form} and the pullback of $\phi$ have
the same divisor on the normal projective variety $R$.  Their
ratio has neither zeros nor poles and is a global unit.  Since
$R$ is geometrically integral, that unit belongs to $K_1^*$.
\end{proof}

For later use, put $L=K_1(R)$ and let $K_2\supset K_1$ be a finite extension Galois
over $K$, and choose a branch over a fixed point of $C$ in the
corresponding tower of curves.  The inertia groups are cyclic and
the upper one surjects onto the lower one.  Choose generators of
these inertia groups, viewed as elements
$\tau_2\in\operatorname{Gal}(K_2/K)$ and
$\tau_1\in\operatorname{Gal}(K_1/K)$ with
$\tau_2|_{K_1}=\tau_1$, and choose a lift
$g\in\operatorname{Gal}(L/K(V))$ of $\tau_1$.
Since $L/K_1$ is regular, it is linearly disjoint from $K_2/K_1$.
The compatible pair $(g,\tau_2)$ therefore defines a finite-order
automorphism $g_2$ of $LK_2$.  The base change $R_{K_2}$ is
geometrically integral, and $g_2^{s!}$ acts regularly on every
base-changed block by the action of $g^{s!}$ together with
$\tau_2^{s!}$ on its constants.  This supplies the compatible
inertia lift and block actions after the further base changes used
for semistable reduction.
